% answers_short.tex —— 2014 年·陕西卷·理科 · 简短答案速查
% 由 tools/make_short_answers.py 自 answers.tex 的 \gksolhead/\gkanswer 提取，
% 并与 2014/2014_数学真题_陕西理_答案要点.md 逐题比对后生成，勿手改（改源 MD 后重跑）。

\begin{gknotice}
\textbf{使用说明}：本节只给最终答案便于快速核对，完整推导见解析版。答案为编写者独立推导，\textbf{非官方标准答案}。
\end{gknotice}

\gkpart{选择题}
\begin{center}\begin{tabular}{ccccccccccc}
\toprule
题号 & 1 & 2 & 3 & 4 & 5 & 6 & 7 & 8 & 9 & 10 \\
\midrule
答案 & B & B & C & C & D & C & D & B & A & A \\
\bottomrule
\end{tabular}\end{center}

\gkpart{填空题}
\begin{itemize}[leftmargin=2.6em]
  \item[\textbf{11.}] $\sqrt{10}$。$4^a=2$ 得 $a=\dfrac12$，故 $x=10^{1/2}=\sqrt{10}$。
  \item[\textbf{12.}] $x^2+(y-1)^2=1$。圆心 $(1,0)$ 关于 $y=x$ 的对称点为 $(0,1)$。
  \item[\textbf{13.}] $\dfrac12$。由 $\boldsymbol a\parallel\boldsymbol b$ 得 $\sin2\theta=\cos^2\theta$，即 $2\sin\theta\cos\theta=\cos^2\theta$，故 $\tan\theta=\dfrac12$。
  \item[\textbf{14.}] $F+V-E=2$。表中三行均满足 $5+6-9=6+6-10=6+8-12=2$。
  \item[\textbf{15.}] A 为 $\sqrt5$；B 为 $3$；C 为 $1$。A 用柯西不等式；B 由 $AE=AC\cos A$ 得 $\cos A=\dfrac12$，再用余弦定理；C 化极坐标方程为 $-x+\sqrt3y=2$ 后求点到直线距离。
\end{itemize}

\gkpart{解答题}
\begin{itemize}[leftmargin=2.6em]
  \item[\textbf{16.}] (1) 见解析版；(2) $\dfrac12$。由 $b^2=ac$ 得 $\cos B=\dfrac{a^2+c^2-ac}{2ac}\geqslant\dfrac12$
  \item[\textbf{17.}] (1) 见解析版；(2) $\dfrac{\sqrt{10}}5$。由三视图得 $DA=1$，$DB=DC=2$ 且两两垂直，平面 $EFGH$ 的法向量为 $(1,1,0)$，与 $\overrightarrow{AB}=(2,0,-1)$ 求线面角
  \item[\textbf{18.}] (1) $2\sqrt2$。$P$ 为 $\triangle ABC$ 的重心 $(2,2)$，故 $|\overrightarrow{OP}|=2\sqrt2$；(2) $m-n=y-x$，最大值为 $1$。由 $x=m+2n$，$y=2m+n$ 解出 $m-n=y-x$，在三角形三个顶点处比较得最大值 $1$
  \item[\textbf{19.}] (1) $P(X=800)=0.2$，$P(X=2000)=0.5$，$P(X=4000)=0.3$。利润 $X=$ 产量 $\times$ 价格 $-1000$，按独立事件相乘后合并；(2) $0.896$。单季利润不少于 $2000$ 元的概率为 $0.8$，三季中至少两季的概率为 $\mathrm C_3^2(0.8)^2(0.2)+(0.8)^3$
  \item[\textbf{20.}] (1) $a=2$，$b=1$。公共点在 $y=0$ 上，由 $C_2$ 得 $x=\pm1$，故 $b=1$；再由 $e=\dfrac{\sqrt3}{2}$ 得 $a=2$；(2) $8x+3y-8=0$。设 $l:y=k(x-1)$，求出 $P$，$Q$ 坐标，由 $\overrightarrow{AP}\cdot\overrightarrow{AQ}=0$ 解得 $k=-\dfrac83$
  \item[\textbf{21.}] (1) $g_n(x)=\dfrac{x}{1+nx}$。由 $g(x)=\dfrac{x}{1+x}$ 迭代并用归纳法即得；(2) $(-\infty,1]$。$x>0$ 时 $a\leqslant\dfrac{(1+x)\ln(1+x)}x$，该式在 $(0,+\infty)$ 上递增且趋于 $1$；(3) 见解析版
\end{itemize}
